\section{Conclusion}
\label{sec:conclusion}

We asked what actually transfers when a pretrained transformer is converted into a
smaller sibling, and answered it in a controlled, single-family setting:
representations align across sizes but parameters do not, dense weight projection
is provably destructive because it breaks the architecture's essential wiring,
and after the best-fit linear operator the weight residuals are indistinguishable
from noise. The value of a conversion therefore lives in initialization, not
in any learnable zero-shot weight correction. Decomposing that initialization into
two independent levers, least-squares compensation (what the network
computes) and variance-preserving rescale (the scale the optimizer sees),
lets us state the practical rule precisely: compensation is a token-efficient,
low-budget win that stacks with rescale to dominate subcloning at small budgets on
both a width- and a depth-reduced pair, ties it once the budget is large enough to
converge, and, stacked, over-corrects only at the largest donor scale, where the
compensation solve is ill-conditioned. For a practitioner spinning up a new size,
the takeaway is concrete: initialize by structure-respecting selection, compensate
and rescale on disjoint paths for a strong low-budget start, and fall back to the
single robust rescale lever at large scale until a dimension-aware regularizer
closes that gap.
